% !TEX TS-program = pdflatex
% Compilar con: pdflatex frontend.tex (2 veces para el índice)

\documentclass[12pt,a4paper]{article}

\usepackage[spanish]{babel}
\usepackage[utf8]{inputenc}
\usepackage[T1]{fontenc}
\usepackage{lmodern}
\usepackage{geometry}
\usepackage{hyperref}
\usepackage{bookmark}
\usepackage{graphicx}
\usepackage{longtable}
\usepackage{array}
\usepackage{xcolor}
\usepackage{listings}
\usepackage{enumitem}
\usepackage{titlesec}

\geometry{margin=2.5cm}

\hypersetup{
  colorlinks=true,
  linkcolor=blue,
  urlcolor=blue,
  citecolor=blue,
  pdftitle={Documentaci\'on Frontend - Plataforma de Aprendizaje C},
  pdfauthor={Proyecto Plataforma C}
}

\definecolor{codebg}{RGB}{248,248,248}
\definecolor{codeborder}{RGB}{230,230,230}
\definecolor{codecomment}{RGB}{90,90,90}

\lstset{
  basicstyle=\ttfamily\small,
  columns=fullflexible,
  breaklines=true,
  frame=single,
  rulecolor=\color{codeborder},
  backgroundcolor=\color{codebg},
  commentstyle=\color{codecomment},
  showstringspaces=false,
  tabsize=2
}

% Algunas instalaciones de MiKTeX no incluyen JavaScript en listings por defecto.
% Definimos un lenguaje mínimo para evitar el error de compilación.
\lstdefinelanguage{JavaScript}{
  morekeywords={break,case,catch,class,const,continue,debugger,default,delete,do,else,export,extends,finally,for,function,if,import,in,instanceof,let,new,return,super,switch,this,throw,try,typeof,var,void,while,with,yield,async,await,from},
  sensitive=true,
  morecomment=[l]{//},
  morecomment=[s]{/*}{*/},
  morestring=[b]',
  morestring=[b]"
}

\title{Documentaci\'on Frontend (React)\\Plataforma de Aprendizaje de Lenguaje C}
\author{Repositorio: plataforma-aprendizaje-c}
\date{Enero 2026}

\begin{document}
\maketitle
\tableofcontents
\newpage

\section{Resumen ejecutivo}
Este documento describe el \textbf{frontend} de la plataforma: arquitectura, ruteo, sesi\'on/autenticaci\'on, componentes y estilos.

El objetivo del frontend es:
\begin{itemize}[leftmargin=*]
  \item Proveer una experiencia SPA moderna construida en \textbf{React}.
  \item Reutilizar componentes comunes (\textbf{Header de usuario} y \textbf{Footer}) en toda la aplicaci\'on.
  \item Integrar el contenido existente (m\'as de 100 HTML) sin romper funcionalidades interactivas como las \textbf{preguntas}.
\end{itemize}

\section{Estructura del proyecto (frontend)}
La parte React se encuentra en:\\
\texttt{Proyecto de info/react-spa/}

Archivos principales:
\begin{itemize}[leftmargin=*]
  \item \texttt{src/main.jsx}: entrada de React, monta Router y SessionProvider.
  \item \texttt{src/App.jsx}: definici\'on de rutas.
  \item \texttt{src/index.css}: estilos globales (layout, hub, viewer, etc.).
  \item \texttt{src/components/}: componentes reutilizables (Shell, HeaderUsuario, Footer, etc.).
  \item \texttt{src/pages/}: pantallas (Landing, Login, Cursos, Viewer, etc.).
  \item \texttt{src/context/SessionContext.jsx}: estado global de sesi\'on (usuario logueado).
\end{itemize}

\section{Tecnolog\'ias y librer\'ias}
\begin{itemize}[leftmargin=*]
  \item \textbf{React}: UI declarativa.
  \item \textbf{React Router}: ruteo SPA con \texttt{BrowserRouter}.
  \item \textbf{Vite}: bundler y tooling de desarrollo/build.
\end{itemize}

\subsection{Base path (deploy bajo /react)}
El frontend est\'a construido para vivir bajo el prefijo \texttt{/react}.

\begin{itemize}[leftmargin=*]
  \item Vite: \texttt{base: '/react/'} (ver \texttt{vite.config.js}).
  \item React Router: \texttt{basename="/react"} (ver \texttt{src/main.jsx}).
\end{itemize}

\section{C\'omo se sirve el frontend}
Aunque este documento es del frontend, es importante entender la integraci\'on:

\begin{itemize}[leftmargin=*]
  \item Express sirve el build est\'atico de React desde \texttt{react-spa/dist} en la ruta \texttt{/react}.
  \item Se configura un \textbf{fallback} para que \texttt{/react/...} devuelva \texttt{index.html} y React Router resuelva la ruta.
  \item Rutas antiguas HTML (por ejemplo \texttt{/login}, \texttt{/registro}, \texttt{/cursos-selec.html}) se redirigen hacia rutas React equivalentes.
\end{itemize}

\section{Ruteo del frontend (React Router)}
Las rutas del frontend se definen en \texttt{src/App.jsx}.

\subsection{Rutas p\'ublicas}
\begin{longtable}{|>{\raggedright\arraybackslash}p{3.5cm}|>{\raggedright\arraybackslash}p{10cm}|}
\hline
\textbf{Ruta} & \textbf{Pantalla} \\
\hline
\texttt{/} & Landing (entrada p\'ublica) \\
\hline
\texttt{/login} & Login (POST /login) \\
\hline
\texttt{/registro} & Registro (POST /register) \\
\hline
\texttt{/forgot-password} & Recuperar (POST /forgot-password) \\
\hline
\texttt{/restore-password} & Restablecer (POST /restore-password, usa query \texttt{token}) \\
\hline
\end{longtable}

\subsection{Rutas protegidas}
Estas rutas requieren sesi\'on (controladas por \texttt{RequireAuth}):
\begin{longtable}{|>{\raggedright\arraybackslash}p{3.5cm}|>{\raggedright\arraybackslash}p{10cm}|}
\hline
\textbf{Ruta} & \textbf{Pantalla} \\
\hline
\texttt{/cursos} & Hub de cursos: lista de m\'odulos y lecciones \\
\hline
\texttt{/m/:modIndex/:fileIndex} & Viewer de contenido con anterior/siguiente \\
\hline
\texttt{/inicio} & Inicio (placeholder/migraci\'on m\'inima) \\
\hline
\texttt{/admin} & Admin (placeholder/migraci\'on m\'inima) \\
\hline
\texttt{/configuracion} & Configuraci\'on usuario (placeholder/migraci\'on m\'inima) \\
\hline
\end{longtable}

\section{Sesi\'on y autenticaci\'on}
El frontend usa cookies de sesi\'on del backend (\texttt{express-session}).

\subsection{SessionProvider}
\texttt{src/context/SessionContext.jsx} define un contexto global que:
\begin{itemize}[leftmargin=*]
  \item Llama a \texttt{GET /search-user} para obtener el usuario.
  \item Mantiene \texttt{status}: \texttt{loading | authenticated | guest}.
  \item Expone \texttt{refresh()} para revalidar sesi\'on.
  \item Expone \texttt{logout()} (POST /logout + refresh).
\end{itemize}

Puntos clave:
\begin{itemize}[leftmargin=*]
  \item Todas las llamadas usan \texttt{credentials: 'include'} para enviar cookies.
  \item Cuando el usuario inicia sesi\'on desde Login, se ejecuta \texttt{refresh()}.
\end{itemize}

\subsection{RequireAuth}
\texttt{src/components/RequireAuth.jsx} bloquea rutas protegidas:
\begin{itemize}[leftmargin=*]
  \item Si status es \texttt{loading}: muestra un estado de carga.
  \item Si es guest: redirige a \texttt{/login?from=...}.
  \item Si es authenticated: renderiza el contenido.
\end{itemize}

\section{Layout y componentes repetidos}
\subsection{Shell}
\texttt{src/components/Shell.jsx} es el layout principal:
\begin{itemize}[leftmargin=*]
  \item Inserta \textbf{HeaderUsuario} al inicio.
  \item Renderiza el \textbf{contenido de la p\'agina}.
  \item Inserta \textbf{Footer} al final.
\end{itemize}

\subsection{HeaderUsuario}
\texttt{src/components/HeaderUsuario.jsx} muestra:
\begin{itemize}[leftmargin=*]
  \item Nombre de usuario cuando est\'a autenticado.
  \item Bot\'on de cerrar sesi\'on que llama \texttt{logout()} y vuelve al inicio.
  \item Links de navegaci\'on (Cursos, Configuraci\'on).
\end{itemize}

\subsection{Footer}
\texttt{src/components/Footer.jsx} contiene estructura de pie de p\'agina (logo C, enlaces r\'apidos, copyright).

\section{Datos: \`Indice de m\'odulos}
El hub de cursos consume un \textbf{\'indice} de m\'odulos generado desde los HTML existentes.

\subsection{Formato}
Se lee desde:\\
\texttt{/js-spa/data/modulosIndex.json}

Estructura esperada:
\begin{lstlisting}[language=JavaScript]
[
  {
    "modulo": "modulo-1-Fundamentos-Introduccion",
    "archivos": [
      { "nombre": "100-Introduccion...html", "ruta": "/Modulos/.../100-...html" }
    ]
  }
]
\end{lstlisting}

\subsection{Generaci\'on del \`indice}
El script \texttt{scripts/generateModulosIndex.js} recorre:\\
	exttt{proyecto/dist/protected\_html/Modulos}

y genera:
\begin{itemize}[leftmargin=*]
  \item \texttt{js-spa/data/modulosIndex.js} (compatibilidad con SPA antiguo)
  \item \texttt{js-spa/data/modulosIndex.json} (consumido por React)
\end{itemize}

\section{P\'agina de Cursos (Hub)}
Implementada en \texttt{src/pages/CursosSelec.jsx}.

Flujo:
\begin{itemize}[leftmargin=*]
  \item Usa \texttt{useModulosIndex()} para cargar el \`indice.
  \item Muestra tarjetas de m\'odulo y links a lecciones.
  \item Incluye barra superior de progreso (TopProgressBar).
\end{itemize}

\subsection{Progreso del usuario}
Actualmente el progreso se obtiene desde \texttt{useUserProgress()} que llama a una API futura \texttt{/api/progress}.
Si no existe, se muestra 0\% sin romper la UI.

\section{Viewer de contenido y por qu\'e usa iframe}
La mayor parte del contenido son documentos HTML completos (con head, scripts y CSS propio). Inyectarlos como \texttt{innerHTML} en una SPA suele romper:
\begin{itemize}[leftmargin=*]
  \item ejecuci\'on de scripts del documento,
  \item carga de estilos declarados en \texttt{<head>},
  \item selectores globales que asumen documento completo,
  \item funcionalidades interactivas (por ejemplo: p\'aginas de preguntas).
\end{itemize}

Por eso \texttt{src/pages/ModuloViewer.jsx} renderiza:
\begin{itemize}[leftmargin=*]
  \item \texttt{<iframe src={file.ruta} .../>}
  \item Navegaci\'on \textbf{Anterior/Siguiente} calculada desde un arreglo plano (flat) de lecciones.
\end{itemize}

Esto preserva el comportamiento original del HTML.

\section{CSS y estilos}
\subsection{D\'onde se editan estilos del React}
El archivo principal es \texttt{src/index.css}. All\'i est\'an:
\begin{itemize}[leftmargin=*]
  \item Estilos del hub (\texttt{.hub-*}).
  \item Estilos del viewer (\texttt{.viewer-*}).
  \item Estilos de navbar cards (\texttt{.navbar-*}).
  \item Estilos base (tipograf\'ia, body, links, etc.).
\end{itemize}

\subsection{Importante: estilos dentro del contenido (m\'odulos)}
Como el contenido se carga dentro de un \texttt{iframe}, los estilos del React \textbf{no} afectan el interior del contenido.

Para modificar el look de los m\'odulos, debes editar los HTML/CSS que esos HTML referencian, o bien crear una estrategia com\'un (por ejemplo: enlazar un CSS global en todos los HTML).

\section{Build, despliegue y comandos}
\subsection{Desarrollo}
\begin{itemize}[leftmargin=*]
  \item Backend: \texttt{cd server} y \texttt{npm run dev}
  \item Frontend React (opcional): \texttt{cd react-spa} y \texttt{npm run dev}
\end{itemize}

\subsection{Producci\'on}
\begin{itemize}[leftmargin=*]
  \item \texttt{cd react-spa}
  \item \texttt{npm run build}
  \item Express servir\'a \texttt{react-spa/dist} bajo \texttt{/react}
\end{itemize}

\subsection{Regenerar \`indice de m\'odulos}
\begin{itemize}[leftmargin=*]
  \item \texttt{node scripts/generateModulosIndex.js}
\end{itemize}

\section{Problemas comunes y diagn\'ostico}
\subsection{Me redirige a login aunque estoy logueado}
\begin{itemize}[leftmargin=*]
  \item Verifica que el backend est\'e corriendo en el mismo dominio (\texttt{localhost:8080}).
  \item Revisa en DevTools que la cookie de sesi\'on exista y se env\'ie.
  \item Confirma que \texttt{GET /search-user} responda \texttt{success: true}.
\end{itemize}

\subsection{No aparece la lista de m\'odulos}
\begin{itemize}[leftmargin=*]
  \item Confirma que existe \texttt{js-spa/data/modulosIndex.json}.
  \item Ejecuta el script generador y recarga.
\end{itemize}

\subsection{El contenido en iframe no carga}
\begin{itemize}[leftmargin=*]
  \item Revisa que la ruta \texttt{/Modulos/...} exista y el backend la sirva.
  \item Si el backend protege la ruta, aseg\'urate de estar autenticado.
\end{itemize}

\section{Extensiones futuras (sugeridas)}
\begin{itemize}[leftmargin=*]
  \item Implementar \texttt{/api/progress} real (por usuario) para el TopProgressBar.
  \item Migrar gradualmente \texttt{/inicio}, \texttt{/admin}, \texttt{/configuracion} a vistas completas.
  \item Agregar un \texttt{vite proxy} para desarrollar React contra el backend en modo \texttt{npm run dev} sin problemas de cookies.
\end{itemize}

\end{document}
